Der Kopf habe das Volumen $n_\text{St}\,V$ mit $n_\text{St} = \nStO$, der Griff das Volumen $n_\text{H}\,V$ mit $n_\text{H} = \nHoO$. Für Stahl nehmen wir Kohlenstoffstahl, $\rho_\text{St} = \rStO$, für den Griff trockenes Buchenholz, $\rho_\text{H} = \rHoO$ (FoTaBe S.~187). Die mittlere Dichte ist die Gesamtmasse durch das Gesamtvolumen:
\begin{align}
 \bar\rho &= \frac{\rho_\text{St}\,n_\text{St}V+\rho_\text{H}\,n_\text{H}V}{n_\text{St}V+n_\text{H}V} = \rHaF \\
   &= \frac{\nStO\cdot\rStO+\nHoO\cdot\rHoO}{\nStO+\nHoO} \\
   &= \rHag \approx \result{\rHagP{0}{2}}
\end{align}
